% !TeX program = xelatex
\documentclass{beamer}
\usepackage{amsfonts,amsmath,oldgerm,booktabs,tikz}
\usetikzlibrary{shapes.geometric, arrows.meta, positioning}
\usetheme{sintef}
\usepackage{xeCJK}

\usepackage{ragged2e}
\newcommand{\testcolor}[1]{\colorbox{#1}{\textcolor{#1}{test}}~\texttt{#1}}

\usefonttheme[onlymath]{serif}

\titlebackground*{assets/background}

\newcommand{\hrefcol}[2]{\textcolor{cyan}{\href{#1}{#2}}}

\tikzset{
    startstop/.style = {rectangle, rounded corners, minimum width=3cm, minimum height=1cm, text centered, draw=black},
    io/.style = {trapezium, trapezium left angle=70, trapezium right angle=110, minimum width=3cm, minimum height=1cm, text centered, draw=black},
    process/.style = {rectangle, minimum width=3cm, minimum height=1cm, text centered, draw=black},
    decision/.style = {diamond, minimum width=3cm, minimum height=1cm, text centered, draw=black},
    arrow/.style = {thick, ->, >=Stealth}
}

\title{2004年奥运会举办城市的推选过程}
% \subtitle{报告副标题}
\course{数学建模}
\author{陈子健\ 钟敦复}
% \IDnumber{1234567}
\date{2026年6月}

\begin{document}
\maketitle

\section{各种选举方法的介绍}

\begin{frame}
    \begin{itemize}
        \item 简单多数法 \pause 
        \item 单论决胜法 \pause
        \item 系列决胜法 \pause
        \item Coombs法 \pause
        \item Borda计数法
    \end{itemize}
\end{frame}

\begin{frame}{系列决胜法}
    \begin{center}
        \begin{tikzpicture}[scale=0.5, transform shape] 
            \node (start) [startstop] {开始};
            \node (pro1) [process, below=of start] {计算每个候选人的得票数};
            \node (dec1) [decision, below=of pro1] {是否仅剩下一名候选人？};
            \node (pro2) [process, right=of dec1] {删除位次第一得票数最少的候选人};
            
            \node (pro3) [process, below=of dec1] {剩下候选人即为获胜者};
            \node (stop) [startstop, below=of pro3] {结束};

            \draw [arrow] (start) -- (pro1);
            \draw [arrow] (pro1) -- (dec1);
            \draw [arrow] (dec1) -- node[anchor=east] {是} (pro3);
            \draw [arrow] (dec1) -- node[anchor=south] {否} (pro2);
            \draw [arrow] (pro2) |- (pro1);
            \draw [arrow] (pro3) -- (stop);

        \end{tikzpicture}
    \end{center}
\end{frame}

\begin{frame}{Coombs法}
    \begin{center}
        \begin{tikzpicture}[scale=0.5, transform shape] 
            \node (start) [startstop] {开始};
            \node (pro1) [process, below=of start] {计算每个候选人的得票数};
            \node (dec1) [decision, below=of pro1] {是否仅剩下一名候选人？};
            \node (pro2) [process, right=of dec1] {删除位次倒数第一得票数最多的候选人};
            
            \node (pro3) [process, below=of dec1] {剩下候选人即为获胜者};
            \node (stop) [startstop, below=of pro3] {结束};

            \draw [arrow] (start) -- (pro1);
            \draw [arrow] (pro1) -- (dec1);
            \draw [arrow] (dec1) -- node[anchor=east] {是} (pro3);
            \draw [arrow] (dec1) -- node[anchor=south] {否} (pro2);
            \draw [arrow] (pro2) |- (pro1);
            \draw [arrow] (pro3) -- (stop);

        \end{tikzpicture}
    \end{center}
\end{frame}

\section{2004 年奥运会举办城市的推选过程}

\begin{frame}
    \begin{table}{2004 年奥运会举办城市各轮投票结果}
        \begin{tabular}{cccccc}
            \toprule
            城市 & 雅典 & 布宜诺斯艾利斯 & 开普敦 & 罗马 & 斯德哥尔摩 \\
            \midrule 
            \pause 第1轮票数 & 32 & 16 & 16 & 23 & 20 \\
            \pause 第2轮票数 & 38 &    & 22 & 28 & 19 \\
            \pause 第3轮票数 & 52 &    & 20 & 35 &    \\
            \pause 第4轮票数 & 66 &    &    & 41 &    \\
            \bottomrule
        \end{tabular}
    \end{table}
\end{frame}

\section{采用系列决胜法在推选过程中可能出现的情况}

\begin{frame}
    \begin{overprint}
        \onslide<2-3> 尽管开普敦在第1轮投票中是得票最少的城市之一，但并不是在淘汰布宜诺斯艾利斯后下一个被淘汰的。
        \onslide<4-5> 如果第2轮投票中投给斯德哥尔席的19票在第3轮投票中都投给开普敦，那么开普敦将得到22+19=41票，第3轮淘汰的将是罗马，而开普敦将在最后一轮投票中与雅典对决。二者之中哪个会获胜？
        \onslide<6-> 如果投给罗马的35票在最后一轮投票中都投给开普敦，那么开普敦将是最终的获胜者。
    \end{overprint}
    \begin{table}{2004 年奥运会举办城市各轮投票结果}
        \begin{tabular}{cccccc}
            \toprule
            城市 & 雅典 & 布宜诺斯艾利斯 & 开普敦 & 罗马 & 斯德哥尔摩 \\
            \midrule 
            \onslide<1-> 第1轮票数 & 32 & 16 & 16 & 23 & 20 \\
            \onslide<3-> 第2轮票数 & 38 &    & 22 & 28 & 19 \\
            \onslide<5-> 第3轮票数 & 38 &    & 41 & 28 &    \\
            \onslide<7-> 第4轮票数 & 38 &    & 69 &    &    \\
            \bottomrule
        \end{tabular}
    \end{table}
        \onslide<8> 这就是说，采用系列决胜法在最初一轮投票中差点被淘汰的城市也可能最终胜！
\end{frame}

\section{采用系列决胜法可能违反单调性准则}

\begin{frame}
    \begin{table}{2004年奥运会举办城市各轮投票结果 \\ （假定第1轮投给开普敦的票中有一票转投给雅典）}
        \begin{tabular}{cccccc}
            \toprule
            城市 & 雅典 & 布宜诺斯艾利斯 & 开普敦 & 罗马 & 斯德哥尔摩 \\
            \midrule 
            \pause 第1轮票数 & 33 & 16 & 15 & 23 & 20 \\
            \pause 第2轮票数 & 33 & 31 &    & 23 & 20 \\
            \pause 第3轮票数 & 33 & 51 &    & 23 &    \\
            \bottomrule
        \end{tabular}
    \end{table}
\end{frame}

\begin{frame}
	我们看到，在用系列决胜法的过程中，第1轮投票雅典位居第1，有利于雅
	典的一点点改变（投给开普敦的票中有一票转投给雅典），在最终排序中却可以使雅典
	落到第2位。如果真的如此（这完全是可能的），单调性准则就不成立了。
\end{frame}

\backmatter
\end{document}
